\documentclass[11pt]{article}
\usepackage[margin=0.85in]{geometry}
\usepackage{amsmath,booktabs,xurl,hyperref}
\usepackage[T1]{fontenc}
\setlength{\emergencystretch}{2em}
\setlength{\tabcolsep}{4pt}
\title{Why the CLF Does Not Return the Ego to Nominal Cruise\\After the Encounter}
\author{Pan Dynamics Collision Avoidance Research}
\date{October 1, 2026}
\begin{document}
\maketitle

\section{Question and answer}
Once the target is beyond the 30-m encounter range, the controller has no
collision rows. Its only link to the nominal behaviour, the given path at the
reference speed, is the CLF stage. In the 5-cm campaign
(\path{COLLISION_MARGIN_5CM_20261001.tex}), 12 of the 14 cases reach this phase.
Three return to nominal cruise within 80~s, at 43.85--62.65~s; nine do not. This
report asks why the CLF does not dissipate the ego state in this phase.

\paragraph{Answer.} The CLF is a decrease condition for one 50-ms hold, at a
fixed 1\% rate, on the LQR quadratic of the cruise trim, and it is the only
term in the problem that refers to the path. Nothing gives the remaining 67--75
planned inputs a recovery purpose, and the real-time iteration (RTI) centers
the next frame's trust region on that unpurposed plan. Four measured
mechanisms follow:
\begin{enumerate}
\item \emph{Trust-region trap.} It governs the closed loop in all nine failures
and by itself explains the five near the path. The later planned
inputs cancel a median 78--83\% of each first-step steering correction. The
shifted plan then becomes the next anchor, which leaves a net steering
authority of 0.011--0.016~rad about the straight trim. A 200-m fixture curve
alone needs 0.016--0.017~rad. Both curved head-on runs therefore reach the path
and leave it again. Far from the path, this authority cannot turn the vehicle
around, and there the one-step condition also fails (item 3). In the braking channel the same loop produces a speed cycle with a
period of about 5~s, which keeps all three 15-m/s straight cases from
confirming recovery.
\item \emph{No recursive decrease.} In the laterally trapped frames, no later
hold of the plan decreases $V$ by 1\%; in the braking traps, 30--36\% do. In 10
of 11 trapped cases the median planned endpoint has a larger $V$ than the
current state. In trapped frames the shifted plan never meets the next frame's
decrease condition.
\item \emph{The one-step condition is not a valid CLF condition outside a
neighbourhood of the path.} The 1\% one-step decrease is unattainable for
headings that point away from the path by $10^\circ$ or more. At zero heading
error it is unattainable beyond 5--7~m of lateral offset, and at every heading beyond about
$2v/\lambda$ (80~m at 8~m/s). Where it is unattainable, the slack turns the
stage into a greedy minimizer of the next-hold value. Far from the path that
minimizer brakes, slides and leaves the tire model's domain: all four idealized
one-step policies, which run without RTI, fail in 4 of 6 far starts.
\item \emph{Slow rate where it works.} Where the decrease is attainable, the
solver treats every input that achieves 1\% as equally good. The realized
1.2--2.4\% per hold costs 14--23~s even without the trap.
\end{enumerate}
The two crossing cases never reach this phase. In both, the ego escorts the
target, so the encounter never ends.

\section{What the controller asks of the CLF in this phase}
The CLF is $V=\|Fe\|^2=e^TPe$ with
$e=[e_y,\,e_\psi,\,e_v,\,v_y-v_y^*,\,r-r^*]$ in the path frame. $P$ is the
discrete LQR matrix of the cruise trim of the path's curvature, and $e_\psi$ is
wrapped to $(-\pi,\pi]$. Each frame solves two conic programs on one affine
model of the plan, which has 68 holds (3.4~s) at 8~m/s and 76 holds (3.8~s) at
15~m/s.
\begin{itemize}
\item \emph{Stage 1} minimizes the collision slacks. Beyond the encounter range
there are no collision rows, so this optimum is 0.
\item \emph{Stage 2} minimizes $\rho$ subject to
$V(x_1)\le0.99\,V(x_0)+\rho$, where $x_1$ is the predicted state after the
\emph{first} hold.
\end{itemize}
The remaining constraints are the same in both stages:
\begin{itemize}
\item the affine dynamics;
\item the trust boxes $|u_i-\bar u_i|\le(0.075,\,0.125)$ and
$|x_i-\bar x_i|\le(2.5,2.5,0.25,2.5,1.5,0.75)$;
\item the braking rate rows;
\item the physical state limits;
\item a second-order-cone endpoint condition on the intrinsic state
$(v_x,v_y,r)$ and the last input.
\end{itemize}
The endpoint core is pose-free and is built for the \emph{straight} trim
(\texttt{terminalContinuation.build(cfg,0)}), on curved roads too. The anchor
$(\bar x,\bar u)$ is the previous plan shifted by one hold, with the straight
trim input appended, rolled out from the measured state. The lane-feedback flow
seed is used only at the first frame or when the warm start fails. Thus, without a target,
every term that refers to the path acts on one hold.

\paragraph{Data.} Instrumented 80-s replays of all 14 cases use a frozen copy of
the controller at commit \path{88ed3c2}. They use the same fixtures, ODE45 plant
and exact state as the campaign. Each frame stores the issued input, the
returned plan and affine states, the anchor, and the CLF values. The replays
reproduce the campaign's recovery times. All analyses are offline, and no
production file changed.

\section{Where the no-risk phase goes}
A frame is \emph{free} when its encounter-exit index is 0, that is, the target
is beyond range at the current state. Free frames are classified as
\begin{itemize}
\item $\rho=0$: the 1\% decrease is met;
\item steering edge: $\rho>0$ and the issued steering is on its box edge;
\item braking edge: $\rho>0$, steering interior, braking on its edge.
\end{itemize}
\begin{center}\footnotesize\setlength{\tabcolsep}{3pt}
\begin{tabular}{rlrrrrrrrrr}
\toprule
 & & Start & $e_y$ at & Closest & Final & Final & \multicolumn{3}{c}{Free frames} & Recovered\\
Speed & Scenario & (s) & start (m) & $|e_y|$ (m) & $e_y$ (m) & $e_v$ (m/s) & $\rho=0$ & steer edge & brake edge & (s)\\
\midrule
8 & headOn & 4.85 & 8.00 & 0.000 & 0.000 & 0.000 & 1007 & 496 & 0 & 57.10\\
8 & acceleratingHeadOn & 4.10 & 7.93 & 0.000 & $-0.001$ & $-0.001$ & 933 & 585 & 0 & 62.65\\
8 & brakingLead & 21.05 & $-7.33$ & 0.000 & 0.004 & $-0.003$ & 1177 & 2 & 0 & 43.85\\
8 & turningCrossing & 16.10 & $-17.52$ & 17.95 & $-97.54$ & 0.455 & 0 & 1272 & 0 & ---\\
8 & curvedHeadOn & 5.05 & 7.68 & 0.016 & $-72.26$ & $-0.857$ & 148 & 1351 & 0 & ---\\
8 & curvedCrossing & 29.80 & 38.93 & 25.44 & 25.44 & $-0.388$ & 89 & 915 & 0 & ---\\
15 & headOn & 2.95 & 4.23 & 0.000 & $-0.010$ & $-0.339$ & 1125 & 0 & 416 & ---\\
15 & acceleratingHeadOn & 2.70 & 4.08 & 0.000 & $-0.007$ & $-0.242$ & 1139 & 0 & 407 & ---\\
15 & brakingLead & 33.35 & $-4.83$ & 0.000 & 0.005 & $-0.226$ & 710 & 33 & 190 & ---\\
15 & turningCrossing & 13.75 & 37.90 & 38.16 & 521.89 & $-1.773$ & 0 & 1254 & 0 & ---\\
15 & curvedHeadOn & 3.90 & 4.02 & 0.007 & $-26.95$ & $-1.787$ & 98 & 1355 & 0 & ---\\
15 & curvedCrossing & 25.70 & $-40.50$ & 40.84 & $-139.65$ & $-1.518$ & 0 & 1086 & 0 & ---\\
\bottomrule
\end{tabular}\end{center}
``Start'' is the first free frame at or after 2~s. The counts cover the whole
replay. 15-m/s turningCrossing also has 71 interior frames.
Final errors are at recovery or at 80~s. The crossing cases at both speeds have
no free frames, and their final $e_y$ is $-621.0$ and $-629.0$~m. In the six
cases that fail laterally, the issued steering is on its box edge in 90--100\%
of free frames, and $V$ grows by a median 0.11--0.35\% per hold in those frames.
In the three 15-m/s straight cases, braking alone is on its edge in 20--27\% of
free frames.

\section{Mechanism 1: the trust-region trap}
\subsection{The plan cancels the correction, and the shift keeps it}
In each steering-edge frame, the first correction
$\delta_0-\bar\delta_0$ is $\pm0.075$~rad. The table compares it with the
corrections of all later inputs, and reports the \emph{residual}. The residual
is the issued steering minus the straight trim, signed by the first correction.
The straight trim is the input that every shifted plan appends at its tail.
\begin{center}\footnotesize
\begin{tabular}{rlrrrrr}
\toprule
 & & Edge & Later sum & Median share & Residual (rad) & Median $V$ ratio\\
Speed & Scenario & frames & opposite & cancelled & median [p10, p90] & per hold\\
\midrule
8 & headOn & 496 & 100\% & 82\% & 0.0131 [0.0109, 0.0154] & 1.0027\\
8 & acceleratingHeadOn & 585 & 100\% & 82\% & 0.0139 [0.0111, 0.0155] & 1.0026\\
8 & turningCrossing & 1272 & 99.9\% & 80\% & 0.0147 [0.0133, 0.0161] & 1.0011\\
8 & curvedHeadOn & 1351 & 100\% & 81\% & 0.0144 [0.0118, 0.0159] & 1.0025\\
8 & curvedCrossing & 915 & 100\% & 80\% & 0.0148 [0.0135, 0.0166] & 1.0011\\
15 & brakingLead & 33 & 100\% & 78\% & 0.0112 [0.0074, 0.0158] & 1.0011\\
15 & turningCrossing & 1254 & 100\% & 83\% & 0.0131 [0.0123, 0.0137] & 1.0035\\
15 & curvedHeadOn & 1355 & 100\% & 79\% & 0.0158 [0.0140, 0.0177] & 1.0030\\
15 & curvedCrossing & 1086 & 100\% & 82\% & 0.0134 [0.0120, 0.0154] & 1.0014\\
\bottomrule
\end{tabular}\end{center}
The 2 edge frames of 8-m/s brakingLead are omitted. In 99.8--100\% of the edge
frames, the second planned correction also opposes the first; its median size
is only 0.003--0.006~rad. The cancellation is spread over about a dozen later
inputs.

This repeats every frame:
\begin{enumerate}
\item The CLF pushes the first input to the edge that turns the vehicle back.
\item Only the first input enters $\rho$, so the later inputs have no objective.
\item The interior-point solution keeps the later state deviations near the
anchor. In \path{RTI_FAILURE_ROOT_CAUSES_20261001.tex}, fixing every later
input at its anchor changed the minimized slack by at most 0.35\%. Removing the
state box removed most of the cancellation, although that box was only 15\%
used at the solution.
\item The shift makes the counter-corrected second input the next anchor's
first input, so the next box again excludes the turning action.
\end{enumerate}

In this stationary regime, each input enters the plan as the straight trim,
collects about $-0.06$~rad of cancellation while it moves to the head, and then receives the $+0.075$~rad edge
correction. The issued steering is therefore pinned about 0.014~rad from the
straight trim, on the side that turns back. The CLF stage still minimizes
$\rho$, but its effective feedback is now a nearly constant offset of about
0.014~rad.

\subsection{What the residual does in each failing case}
\paragraph{Curved head-on: the ego reaches the path and leaves it.} The fixture
curve has radius 200~m. Its trim steering is 0.0157~rad at 8~m/s and 0.0169~rad
at 15~m/s, and the residual is 0.0144 and 0.0158~rad. From 20 to 80~s, the issued
steering is a median 0.0011~rad \emph{below} the curved trim in both cases. The
anchor's first steering sits 0.076~rad below the curved trim. All the
remaining authority is therefore spent on following the curve, and none is
left for correcting the heading.
\begin{itemize}
\item At 8~m/s, the ego crosses the path at 20~s with $e_\psi=-0.32$~rad. The
heading error stays between $-0.20$ and $-0.28$~rad until 50~s, the ego drifts
outward at a mean 1.2~m/s, and $e_y$ reaches $-72.3$~m at 80~s.
\item At 15~m/s, the ego is on the path at 7~s. It drifts to $-29$~m at 40~s
and then wanders between $-12$ and $-29$~m.
\end{itemize}
The closest approaches were 16 and 7~mm.

\paragraph{Far from the path: no turnaround.} In turningCrossing and
curvedCrossing at both speeds, the residual corresponds to a turn of radius
roughly 200--260~m. The heading errors are 0.5--2.8~rad and the lateral errors
17--41~m, so recovery needs a turnaround within a few seconds. Three of these
runs end farther from the path than they started: turningCrossing at
$-97.5$~m (8~m/s) and 521.9~m (15~m/s), and 15-m/s curvedCrossing at
$-139.6$~m. 8-m/s curvedCrossing closes only from 38.9 to 25.4~m in 50~s.

\paragraph{Straight 8~m/s: slow recovery.} On a straight path, the 0.013~rad
residual corrects slowly. Recovery takes 496 (headOn) and 585
(acceleratingHeadOn) trapped frames before the error falls into the region where
$\rho=0$ (Section~\ref{sec:rate}).

\paragraph{Straight 15~m/s: a speed cycle.} The lateral errors are inside their
tolerances from 19.2~s (headOn), 18.6~s (acceleratingHeadOn) and 51.0~s
(brakingLead) on. The speed error then cycles between 0 and $-0.33$ to
$-0.40$~m/s. It is below $-0.1$~m/s in 51--55\% of that time, and excursions
recur about every 5~s (intervals of 3.2--6.2~s). In headOn, from 41.25 to 41.85~s, $V_0$ falls from
0.0033 to 0.00015 with $\rho\approx0$. Meanwhile the anchor's first braking ratio
drifts from $-0.031$ to $-0.094$, while the cruise trim is $+0.018$.

At 41.95~s the anchor reaches $-0.114$. The box then admits at most $+0.011$,
below the trim, so $\rho>0$ and the vehicle slows. The anchor at 41.65~s shows
the shape responsible:
\begin{itemize}
\item the braking ratio is about $-0.05$ near the head, rises to $+0.10$ near
hold 60, and returns to the trim at the end;
\item the planned speed dips by 0.71~m/s mid-horizon and returns to the
reference at the endpoint.
\end{itemize}
Nothing in the problem prefers this plan to a flat one. Every time this shape
reaches the head, the issued braking is pinned at the box edge.

\section{Mechanism 2: the plan is not a decrease certificate}
For each free frame, $V$ was evaluated along the returned affine plan and at
the first node of the anchor. The anchor's first node is the shifted previous
plan, applied without correction.
\begin{center}\footnotesize
\begin{tabular}{rlrrrrrr}
\toprule
 & & \multicolumn{3}{c}{Trapped frames} & \multicolumn{3}{c}{$\rho=0$ frames with $V_0\ge1$}\\
 & & Later holds & Plan-end & Anchor & Later holds & Plan-end & Anchor meets\\
Speed & Scenario & with 1\% & $V/V_0$ & $V_1/V_0$ & with 1\% & $V/V_0$ & 1\% decrease\\
\midrule
8 & headOn & 0\% & 1.54 & 1.0049 & 21\% & 0.91 & 49\%\\
8 & acceleratingHeadOn & 0\% & 1.51 & 1.0048 & 49\% & 0.50 & 60\%\\
8 & brakingLead & \multicolumn{3}{c}{2 frames} & 38\% & 0.55 & 55\%\\
8 & turningCrossing & 0\% & 1.12 & 1.0013 & --- & --- & ---\\
8 & curvedHeadOn & 0\% & 1.38 & 1.0033 & 79\% & 0.31 & 79\%\\
8 & curvedCrossing & 0\% & 0.81 & 1.0013 & 100\% & 0.23 & 94\%\\
15 & headOn & 30\% & 31.7 & 1.77 & 0\% & 5.22 & 0\%\\
15 & acceleratingHeadOn & 31\% & 32.1 & 1.77 & 0\% & 5.48 & 0\%\\
15 & brakingLead & 36\% & 18.1 & 1.75 & 0\% & 2.02 & 14\%\\
15 & turningCrossing & 0\% & 1.33 & 1.0036 & --- & --- & ---\\
15 & curvedHeadOn & 0\% & 2.36 & 1.0057 & 10\% & 3.50 & 58\%\\
15 & curvedCrossing & 0\% & 1.25 & 1.0019 & --- & --- & ---\\
\bottomrule
\end{tabular}\end{center}
The columns are medians, except ``anchor meets'', which is a share of frames.
``Trapped'' means steering edge or braking edge. In every case, the shifted
plan meets the next frame's 1\% condition in \emph{none} of the trapped frames.
Even in $\rho=0$ frames, it meets the condition in only 49--60\% of frames in
the straight 8-m/s cases and 0--14\% in the straight 15-m/s cases.
The decrease is therefore produced anew by the first correction each frame. It
is never carried by the plan.

In predictive control, Lyapunov decrease normally comes from a horizon value:
the shifted plan plus a terminal law is a feasible candidate whose cost is lower
by the stage cost. Here the shifted plan neither carries the decrease nor ends
in a set that encodes the path, so decrease at one frame implies nothing about
the next.

\section{Mechanism 3: the one-step condition does not reach}
\subsection{Where the 1\% decrease is attainable}
For straight-lane states $(e_y,e_\psi)$ at the reference speed with $v_y=r=0$,
every input on a $141\times41$ grid was propagated one hold with the
controller's own model (one RK4 step). Steering spans $\pm0.7$~rad and braking
$\pm0.999$. Each result was scored with the controller's $V$. Headings use a
$5^\circ$ grid. Positive headings point toward the path for $e_y<0$.
\begin{center}\footnotesize
\begin{tabular}{rrrrr}
\toprule
 & \multicolumn{2}{c}{8 m/s} & \multicolumn{2}{c}{15 m/s}\\
$|e_y|$ (m) & 1\% attainable for $e_\psi$ & best ratio at $e_\psi=0$ & 1\% attainable for $e_\psi$ & best ratio at $e_\psi=0$\\
\midrule
0.5 & $0^\circ$--$20^\circ$ & 0.951 & $0^\circ$--$10^\circ$ & 0.933\\
2 & $0^\circ$--$60^\circ$ & 0.973 & $0^\circ$--$30^\circ$ & 0.960\\
5 & $0^\circ$--$135^\circ$ & 0.988 & $0^\circ$--$70^\circ$ & 0.982\\
10 & $5^\circ$--$165^\circ$ & 0.994 & $5^\circ$--$145^\circ$ & 0.990\\
20 & $10^\circ$--$160^\circ$ & 0.997 & $5^\circ$--$170^\circ$ & 0.995\\
40 & $25^\circ$--$150^\circ$ & 0.998 & $15^\circ$--$165^\circ$ & 0.998\\
60 & $45^\circ$--$135^\circ$ & 0.999 & $20^\circ$--$160^\circ$ & 0.998\\
100 & none & 0.999 & $40^\circ$--$145^\circ$ & 0.999\\
\bottomrule
\end{tabular}\end{center}
For headings pointing away from the path by $10^\circ$ to $160^\circ$, no
admissible input lowers $V$ at all, at either speed or any tested offset. Every
row also shows an apparent decrease at exactly $180^\circ$. That is the
wrapped heading crossing $\pm\pi$: $V$ drops by
$4\pi P_{12}|e_y|$ without any progress. Three properties of the vehicle
explain the pattern; the quadratic $V$ does not account for them.
\begin{itemize}
\item \emph{Relative degree.} Steering acts on $e_y$ only through
$r\rightarrow\psi\rightarrow e_y$. Within one 50-ms hold, it changes mainly
$v_y$ and $r$, and friction bounds the yaw acceleration. When the heading is
wrong, the decrease over that hold is set by the drift, not by the input. A
parallel offset already fails beyond 5~m (8~m/s) or 7~m (15~m/s).
\item \emph{Bounded approach speed.} Where $V\approx P_{11}e_y^2$, the relative
decrease rate is at most $2v|\sin e_\psi|/|e_y|$. A fixed 1\% per hold requires
$\lambda=-\ln(0.99)/0.05=0.201~\mathrm{s^{-1}}$, so the condition fails beyond
$|e_y|\approx2v|\sin e_\psi|/\lambda$, which is 79.6~m at 8~m/s and 149~m at
15~m/s for $e_\psi=90^\circ$. The map agrees: at 100~m, nothing works at 8~m/s,
and only $40^\circ$--$145^\circ$ works at 15~m/s. A quadratic with a fixed
exponential rate cannot be a global CLF for a vehicle whose approach speed is
bounded.
\item \emph{The quadratic's valley is linear-model geometry.} For fixed $e_y$,
$V$ is minimized at $e_\psi=-(P_{12}/P_{22})e_y$, which is $-0.095\,e_y$ at
8~m/s and $-0.096\,e_y$ at 15~m/s. This preferred heading passes $90^\circ$ at
$|e_y|\approx16.4$~m and $\pi$ at about 33~m. It encodes $\dot e_y=v\,e_\psi$, whereas the vehicle obeys
$\dot e_y=v\sin e_\psi$.
\end{itemize}
The quadratic is fine \emph{locally}. For the hold-discretized linear model with
unconstrained input, the best one-step ratio in the worst direction is 0.959
(8~m/s) and 0.966 (15~m/s), so the 1\% decrease is attainable in every
direction. The failures above come from friction saturation, $\sin e_\psi$ and
the heading wrap, not from the LQR weights.

\subsection{What a one-step CLF does without the trap}
To separate mechanism 3 from mechanism 1, four idealized one-step policies were
run from each case's first free state. They have no RTI, no anchor and no
target.
\begin{itemize}
\item Each hold searches the inputs on a grid. \emph{Full} policies use
$71\times21$ inputs over $\pm0.7$~rad and $\pm0.999$. \emph{Slew} variants use
$31\times21$ inputs within $\pm(0.075,0.125)$ of the previously applied input,
with no steering magnitude bound, as in the controller. An $11\times11$ local
refinement follows.
\item The candidates are scored with the controller's model and $V$.
\item \emph{Greedy} applies $\arg\min V(x_1)$. \emph{Indifferent} applies, among
the inputs with the least slack $\rho$, the one closest to the previous input.
\item The plant is the same model with a 5-ms RK4 mesh. Recovery uses the
campaign tolerances and a 5-s dwell.
\end{itemize}
The crossing cases start at 10~s, with the target ignored.
\begin{center}\footnotesize
\begin{tabular}{rlrrrrrl}
\toprule
 & & $e_y,\,e_\psi$ & \multicolumn{4}{c}{Recovered after (s) or outcome} & \\
Speed & Scenario & at start & greedy & greedy slew & indiff. & indiff. slew & Controller\\
\midrule
15 & headOn & 4.2, $-0.04$ & 14.45 & 7.55 & 27.60 & 22.15 & speed cycle\\
15 & acceleratingHeadOn & 4.1, $-0.04$ & conv.$^a$ & 7.50 & 19.55 & 18.20 & speed cycle\\
15 & brakingLead & $-4.8$, $-0.05$ & conv.$^a$ & 7.70 & conv.$^a$ & 22.25 & speed cycle\\
15 & curvedHeadOn & 4.0, $-0.02$ & 7.55 & 7.55 & 15.75 & 22.60 & $-27.0$~m\\
8 & curvedHeadOn & 7.7, 0.24 & 9.75 & 9.80 & 20.90 & tire$^b$ & $-72.3$~m\\
\midrule
8 & turningCrossing & $-17.5$, $-1.43$ & 14.20 & tire & 25.30 & no input$^c$ & $-97.5$~m\\
15 & curvedCrossing & $-40.5$, $-0.49$ & 17.30 & 16.95 & no input & no input & $-139.6$~m\\
8 & curvedCrossing & 38.9, 2.68 & no input & tire & no input & tire & 25.4~m\\
15 & turningCrossing & 37.9, 2.76 & no input & no input & no input & no input & 521.9~m\\
8 & crossing & $-61.0$, $-1.56$ & no input & no input & no input & no input & escort\\
15 & crossing & $-69.4$, $-1.62$ & no input & no input & no input & no input & escort\\
\bottomrule
\end{tabular}\end{center}
$^a$The errors are inside the tolerances from 2.55--12.55~s on. Input
chattering at the grid resolution drives yaw-rate excursions to 0.0105~rad/s,
against the 0.01~rad/s tolerance. These reset the dwell in 1\% of holds; the
lateral error stays within 12~mm.
$^b$The plant left the tire model's operating domain.
$^c$Every candidate input left the model's domain.

\paragraph{Near the path, the CLF works when it is not trapped.} The one-step
CLF brings all five near-path starts inside the tolerances, under three of the
four policies; the fourth fails once. The errors are inside within 2.5--5~s
(greedy) or 10--17~s (indifferent). The controller recovers none of these five.
Its failures are mechanism 1, the steering trap on the curves and the braking
trap on the straight 15-m/s cases.

\paragraph{Far from the path, even the untrapped CLF fails.} 4 of 6 far starts
fail under every policy. The two successes slide sideways, with $|v_y|$ up to
5.4--9.4~m/s and speed down to 1.9--8.5~m/s.

Decomposing the change in $V$ shows why. In 8-m/s crossing, the start has
$e_y=-61$~m and the heading points away from the path. Greedy applies full lock
and braking ratio $-0.999$ for holds 2--9, yet $V$ still rises by 0.3--1.1\% per
hold:
\begin{itemize}
\item the $e_y$ term adds +2500 to +3800 per hold;
\item turning removes only 350 to 1700 per hold through the heading term.
\end{itemize}
The speed weight $P_{33}=20.5$ is small, so trading speed for slower lateral
divergence is the cheapest one-step reduction. Greedy therefore brakes from
7.95 to 4.85~m/s in 0.4~s. The speed then falls to 1.4~m/s, near the model's
1-m/s floor, where every candidate input leaves the domain. In 15-m/s
turningCrossing, full lock builds $v_y$ to 4.8~m/s within 0.6~s, and the
$v_y$ term rises to +480 per hold.

A myopic minimizer of a quadratic far from its centre is not a stabilizing
policy. The slack therefore cannot rescue the infeasible region.

\section{Mechanism 4: when the decrease is met, it is only 1--2\%}\label{sec:rate}
In $\rho=0$ frames, every input that reaches $0.99\,V_0$ is optimal, and the
interior-point method returns a point inside that set. The realized median
ratio per hold is 0.977--0.988 for $V_0>10$, and 0.90--0.95 for $V_0\le10$.
The three recovered 8-m/s cases split as follows:
\begin{center}\footnotesize
\begin{tabular}{lrrrrr}
\toprule
 & Encounter & Trapped frames & Last trapped frame & $\rho=0$ frames after it & Recovered\\
Scenario & ends (s) & before recovery & (s, $e_y$) & (median ratio) & (s)\\
\midrule
headOn & 4.85 & 496 & 42.90, $-1.03$~m & 283 (0.976) & 57.10\\
acceleratingHeadOn & 4.10 & 585 & 46.35, $-1.69$~m & 325 (0.980) & 62.65\\
brakingLead & 21.05 & 2 & 21.10, $-7.32$~m & 454 (0.982) & 43.85\\
\bottomrule
\end{tabular}\end{center}
The trap accounts for 25--29~s of the two head-on recoveries. The final
$\rho=0$ descent takes another 14--23~s. In the idealized runs, indifferent
policies are two to three times slower than greedy ones from the same starts.

\section{Root cause}
The controller has a Lyapunov \emph{condition}, but no Lyapunov
\emph{mechanism}. The decrease is required of one hold, at a fixed rate, on a
function that is valid only near the path. Nothing in the plan, the terminal
set or the frame-to-frame update carries the decrease forward.
\begin{itemize}
\item Near the path the condition is attainable. However, the RTI trust region
is centered on a plan with no recovery objective, and the solver's centering
turns that plan into a counter-correction. The result is a steering authority
about the size of a 200-m curve's trim, plus a braking wave.
\item Far from the path, the condition itself cannot be met. Its slack-minimizing
surrogate is a destructive greedy policy.
\item The pose-free, straight terminal set does not encode the path, so the end
of the plan exerts no pull back toward it.
\end{itemize}

\paragraph{Implications (not implemented or validated here).}
\begin{enumerate}
\item The whole plan needs a recovery objective consistent with the CLF, such as
a horizon sum of $V$ or a tracking cost, or a multi-hold decrease along the
plan. With such an objective the shifted plan becomes the decrease candidate. A
diagnostic third stage of this kind recovered 9 of 10 runs in the RTI report
(6-mm margin).
In five of those nine runs it touched the target, so it must be paired with
inter-sample safety.
\item The decrease condition must be attainable everywhere. Options are a
guidance-type function with a saturated preferred heading, a decrease bound
that grows sublinearly with the error, and a heading treatment without the
$\pm\pi$ jump.
\item In the no-risk phase, the trust region could be centered on a
recovery-consistent rollout. The existing lane-feedback seed is one example.
\item On curves, the terminal set and the appended tail input should use the
curved trim, not the straight one.
\end{enumerate}

\section{Limits}
\begin{itemize}
\item All closed loops are nominal, with exact state and a known target.
\item The feasibility map covers steady states ($v_y=r=0$, reference speed) on a
grid; intermediate states may differ.
\item The idealized policies use input grids and a 5-ms RK4 plant instead of
ODE45. They are diagnostics of the one-step condition, not proposed
controllers.
\item The near/far split is at $|e_y|=10$~m.
\item The cancellation share and residual are medians over trapped frames.
\item ``Not recovered'' means not within 80~s.
\end{itemize}

\section{Artifacts}
\path{report/CLF_NO_RISK_DISSIPATION_20261001/} contains:
\begin{itemize}
\item the regime, cancellation, plan-decay, trajectory, feasibility-map and
idealized-policy tables;
\item every analysis method and log;
\item reproduction notes;
\item a manifest of source and raw-artifact hashes.
\end{itemize}
The replays (MAT, about 500~MB) and the policy traces remain in
\path{/home/zai/.cache/collisionAvoidance/clf-dissipation-20261001/}.
\end{document}
